\documentclass[10pt,conference]{IEEEtran}

\usepackage{amsmath,amssymb,amsfonts}
\usepackage{algorithmic}
\usepackage{graphicx}
\usepackage{textcomp}
\usepackage{booktabs}
\usepackage{multirow}
\usepackage{microtype}
\usepackage{url}
\usepackage{cite}

\hyphenation{op-tical net-works semi-conduc-tor}

\begin{document}

\title{Advanced Clinical Evidence Intelligence System (CEIS) Using Fine-Tuned LAYA: A Non-Autoregressive Multi-Question Decision Architecture for Healthcare NLP}

\author{
    \IEEEauthorblockN{Shailesh Reddy\IEEEauthorrefmark{1} and Karthikeya\IEEEauthorrefmark{2}}
    \IEEEauthorblockA{Department of Computer Science and Engineering\\
    Indian Institute of Information Technology, Design and Manufacturing (IIITDM), Kurnool\\
    Email: \IEEEauthorrefmark{1}524cs0010@iiitk.ac.in, \IEEEauthorrefmark{2}124cs0021@iiitk.ac.in\\
    Model Repository: \url{https://huggingface.co/reddyshailesh1/ceis} $\cdot$ Code: \url{https://github.com/reddyshailesh1/Ceis}}
}

\maketitle

\begin{abstract}
Conventional healthcare Natural Language Processing (NLP) systems face a critical operational dilemma: standard single-label Natural Language Inference (NLI) classifiers lack the semantic expressiveness to evaluate evidence sufficiency or topic presence, while modern generative Large Language Models (LLMs) introduce stochastic hallucinations, high computational latency, and uncalibrated confidence distributions that threaten patient safety. In this work, we present the Clinical Evidence Intelligence System (CEIS), a medical evidence verification framework extending LAYA's non-autoregressive System-1 decision architecture. By formalizing clinical reasoning into five typed decision primitives---comprising choice, noul (boolean), and continuous score questions---CEIS verifies clinical claims against electronic health records, detects subtle document inconsistencies, quantifies evidence sufficiency, and flags clinical risk for practitioner review in a single deterministic forward pass. We fine-tune the architecture on MedNLI using hardware acceleration on an NVIDIA RTX 6000 Ada (48~GB VRAM) with group-aware premise splitting across MIMIC-III patient IDs, 15\% donor mismatch negative sampling, and temperature scaling ($T = 1.1911$). Head-to-head empirical evaluation demonstrates that CEIS eliminates critical semantic head failures observed in existing baseline checkpoints, achieving a 65.3\% loss reduction on topic mention and a 41.0\% loss reduction on evidence sufficiency while preserving core NLI accuracy (79.08\%). We deploy the complete system with a high-throughput FastAPI REST backend and a clinician-facing Streamlit dashboard featuring span highlighting and pairwise cross-note contradiction matrices.
\end{abstract}

\begin{IEEEkeywords}
Clinical Natural Language Inference, Non-Autoregressive Decisions, Medical Evidence Verification, Contradiction Detection, Model Calibration, LAYA Framework, Healthcare Safety.
\end{IEEEkeywords}

\section{Introduction}
Modern Electronic Health Records (EHRs) contain massive volumes of unstructured narrative text, including physician progress notes, emergency department admission histories, pathology reports, and discharge summaries. Clinicians routinely cross-reference patient histories against current clinical documentation to verify prior diagnoses, medication adherence, and drug allergy profiles. Discrepancies between distinct clinical records---such as contradictory medication listings between admission and transfer, or mismatched diagnostic findings between radiology and discharge summaries---represent a leading cause of preventable medical adverse events~\cite{romanov2018lessons}.

To automate the identification of such inconsistencies, clinical NLP has traditionally framed the task as Natural Language Inference (NLI): given an evidence premise $P$ and a hypothesis claim $H$, the objective is to determine whether $P$ entails, contradicts, or is neutral towards $H$. However, standard single-head classifiers (e.g., ClinicalBERT~\cite{alsentzer2019publicly}, PubMedBERT~\cite{gu2021domain}) output only a 3-way distribution over $\{\text{entailment}, \text{neutral}, \text{contradiction}\}$. When a claim is classified as neutral (unsupported), clinicians cannot determine whether the evidence completely omits the medical topic or discusses it ambiguously without conclusive diagnostic proof. 

Conversely, while autoregressive LLMs can generate open-ended textual justifications, their propensity for hallucination (fabricating non-existent clinical findings), prohibitive token-by-token generation latency, and poor probability calibration severely limit their clinical utility~\cite{guo2017calibration}. 

To overcome these fundamental limitations, we propose the \textbf{Clinical Evidence Intelligence System (CEIS)}, extending the open-source \textbf{LAYA} non-autoregressive decision engine~\cite{laya2026}. CEIS formulates clinical evidence intelligence as a multi-task decision problem over typed primitives, evaluating claim entailment, explicit topic mention, evidence sufficiency, and inconsistency severity simultaneously in a single forward pass.

\textbf{Contributions:}
\begin{enumerate}
    \item \textbf{Typed Decision Architecture:} We design a multi-question reasoning schema combining choice, boolean (\textit{noul}), and continuous score questions evaluated simultaneously in a single forward pass.
    \item \textbf{Leakage-Free Multi-Task Fine-Tuning:} We develop a group-aware data pipeline clustering MIMIC-III premise stems to prevent patient-level data leakage, augmented with 15\% donor premise negative sampling.
    \item \textbf{Empirical Head-to-Head Resolution:} We demonstrate that existing baseline checkpoints suffer from acute semantic incoherence across auxiliary heads (e.g., predicting that definitive diagnostic evidence requires more documentation), which CEIS fine-tuning provably eliminates.
    \item \textbf{Confidence Calibration and Safety Gating:} We implement post-hoc temperature scaling ($T = 1.1911$) on held-out calibration records and establish an automated clinician review gating rule.
    \item \textbf{Open-Source Clinical Deployment:} We release the fine-tuned model on Hugging Face (\url{https://huggingface.co/reddyshailesh1/ceis}) alongside full source code, FastAPI REST endpoints, and an interactive Streamlit clinical dashboard on GitHub (\url{https://github.com/reddyshailesh1/Ceis}).
\end{enumerate}

\section{Related Work}
\subsection{Clinical Natural Language Inference}
Domain-specific NLI benchmarks have driven progress in medical language understanding. Romanov and Shivade~\cite{romanov2018lessons} introduced MedNLI, consisting of sentence pairs annotated by clinicians from MIMIC-III critical care notes~\cite{johnson2016mimic}. Pretrained biomedical transformers, such as ClinicalBERT~\cite{alsentzer2019publicly} and PubMedBERT~\cite{gu2021domain}, established strong classification baselines. PubMedQA~\cite{jin2019pubmedqa} formalized biomedical QA over PubMed abstracts using yes/no/maybe grounding. However, existing clinical NLI systems remain restricted to single-head, single-label paradigms.

\subsection{Non-Autoregressive Decision Architectures}
In contrast to autoregressive LLMs that generate token sequences sequentially, non-autoregressive architectures execute all classification heads in parallel over a shared contextual representation. LAYA~\cite{laya2026} introduces typed decision primitives (Choice, Noul, Score) enabling structured multi-question querying. CEIS extends this framework to multi-task clinical evidence verification, guaranteeing bounded sub-50\,ms inference latency and mathematical protection against hallucinatory text generation.

\section{System Architecture \& Typed Primitives}
Let a clinical input state $\mathbf{x} = (C, E)$ denote a clinical claim $C$ paired with an evidence passage $E$. A bidirectional transformer backbone encodes $\mathbf{x}$ into contextual hidden states $\mathbf{H} \in \mathbb{R}^{L \times d}$. CEIS evaluates five typed decision questions in parallel:

\begin{table}[t]
\centering
\caption{Typed Clinical Decision Primitives in CEIS}
\label{tab:primitives}
\resizebox{\columnwidth}{!}{%
\begin{tabular}{@{}llll@{}}
\toprule
\textbf{Question ID} & \textbf{Type} & \textbf{Output Domain} & \textbf{Clinical Semantics} \\
\midrule
\texttt{verdict} & Choice & $\{\text{sup}, \text{unsup}, \text{contra}\}$ & Primary claim entailment status. \\
\texttt{explicitly\_mentioned} & Noul & $\{\text{true}, \text{false}\}$ & Detects if claim topic is in text. \\
\texttt{needs\_more\_evidence} & Noul & $\{\text{true}, \text{false}\}$ & Flags incomplete documentation. \\
\texttt{evidence\_sufficiency} & Score & $[0.00, 1.00]$ & Quantifies proof completeness. \\
\texttt{inconsistency\_severity} & Score & $[0.00, 1.00]$ & Assesses clinical danger of conflict. \\
\bottomrule
\end{tabular}%
}
\end{table}

\subsection{Choice Primitive: Claim Verdict}
The core claim verification question determines whether the evidence supports, does not support, or contradicts the hypothesis:
\begin{equation}
P(\text{verdict} = c \mid \mathbf{x}) = \frac{\exp(\mathbf{w}_c^\top \mathbf{h}_{\text{pool}})}{\sum_{j=1}^3 \exp(\mathbf{w}_j^\top \mathbf{h}_{\text{pool}})}
\end{equation}
where $\mathbf{h}_{\text{pool}}$ is the sequence summary vector and $\mathbf{w}_c$ is the projection vector for class $c \in \{\text{supported}, \text{unsupported}, \text{contradicted}\}$.

\subsection{Noul (Boolean) Primitives}
In LAYA, boolean questions are typed as \textit{noul} primitives, parameterized by a calibrated scalar logit $z_{\text{noul}} \in \mathbb{R}$:
\begin{equation}
P(Y = \text{true} \mid \mathbf{x}) = \sigma(z_{\text{noul}}) = \frac{1}{1 + \exp(-z_{\text{noul}})}
\end{equation}
CEIS evaluates two noul primitives:
\begin{enumerate}
    \item \texttt{explicitly\_mentioned}: Verifies whether the specific clinical condition, laboratory finding, or anatomical structure is directly cited in the narrative.
    \item \texttt{needs\_more\_evidence}: Ascertains whether additional clinical records are mandatory to reach a definitive medical conclusion.
\end{enumerate}

\subsection{Continuous Score Primitives}
Continuous scores quantify clinical metrics on $[0.0, 1.0]$:
\begin{equation}
\mathcal{S}(\mathbf{x}) = \text{clamp}\left(\mathbf{w}_s^\top \mathbf{h}_{\text{pool}} + b_s, 0.0, 1.0\right)
\end{equation}
\texttt{evidence\_sufficiency} outputs an index representing evidence adequacy (0.90 for explicit entailment, 0.20 for unmentioned). \texttt{inconsistency\_severity} triggers exclusively when $\text{verdict} = \text{contradicted}$, measuring the danger of conflicting clinical statements (e.g., lethal drug allergy conflicts vs minor notation discrepancy).

\section{Dataset Pipeline \& Leakage Prevention}
\subsection{MedNLI Cohort \& Premise Clustering}
We utilize the MedNLI benchmark~\cite{romanov2018lessons}, consisting of 11,232 training pairs, 1,395 validation pairs, and 1,422 test pairs. Because clinicians generated multiple hypothesis claims from each clinical sentence, naive random splitting introduces substantial data leakage: identical patient premises appear across train and test partitions. 

CEIS enforces \textbf{group-aware premise splitting}: every example sharing an identical premise stem is assigned to a single partition cluster $G_k$. All cross-validation and calibration splits maintain strict patient-level isolation.

\subsection{Donor Premise Negative Sampling}
To train the auxiliary heads to recognize irrelevant or mismatched clinical text, we introduce a 15\% donor premise mismatch procedure during training. For each batch, 15\% of examples have their true premise replaced with a premise drawn randomly from a disjoint patient record $G_j \neq G_i$. The ground-truth targets for donor mismatch instances are assigned to:
\begin{equation}
\begin{aligned}
y_{\text{verdict}} &= \text{unsupported}, \quad y_{\text{mentioned}} = \text{false}, \\
y_{\text{needs\_more}} &= \text{true}
\end{aligned}
\end{equation}

\section{Multi-Task Fine-Tuning \& Hardware Acceleration}
\subsection{Multi-Task Loss Formulation}
CEIS jointly optimizes a weighted multi-task cross-entropy objective:
\begin{equation}
\mathcal{L}_{\text{total}}(\theta) = \lambda_v \mathcal{L}_{\text{verdict}} + \lambda_m \mathcal{L}_{\text{mentioned}} + \lambda_n \mathcal{L}_{\text{needs\_more}}
\end{equation}
where $\lambda_v = 0.60$, $\lambda_m = 0.20$, and $\lambda_n = 0.20$. Label smoothing ($\alpha = 0.05$) is applied to prevent overconfident peak logit distributions. Silver labels for boolean heads are synthesized deterministically from entailment status and donor mismatch markers.

\subsection{Remote Hardware Acceleration}
Fine-tuning was conducted on a remote compute server configured with:
\begin{itemize}
    \item \textbf{GPU:} 1 $\times$ NVIDIA RTX 6000 Ada Generation (49,140\,MiB / 48\,GB VRAM, 18,176 CUDA cores, 568 Tensor cores)
    \item \textbf{CPU:} 2 $\times$ Intel Xeon Platinum 8480+ (112 physical cores, 224 threads)
    \item \textbf{System Memory:} 123\,GiB RAM
    \item \textbf{Software Stack:} Linux 7.0 (Ubuntu 26.04 LTS), PyTorch 2.13.0, CUDA 13.0, AMP \texttt{bfloat16}
\end{itemize}

With effective batch size 64 (batch size 32 with 2 gradient accumulation steps) and AdamW optimization ($\text{lr}_{\text{enc}} = 2 \times 10^{-5}$, $\text{lr}_{\text{head}} = 10^{-4}$), each 339-step epoch executed in approximately 1.4 minutes. Training converged in \textbf{4.5 minutes} across 3 epochs under early stopping.

\begin{table}[t]
\centering
\caption{Multi-Task Loss Convergence Across Training Epochs}
\label{tab:convergence}
\resizebox{\columnwidth}{!}{%
\begin{tabular}{@{}lccccc@{}}
\toprule
\textbf{Training Phase} & \textbf{Verdict} & \textbf{Mentioned} & \textbf{NeedsMore} & \textbf{Combined} & \textbf{Val Acc} \\
\midrule
Pre-Training Baseline   & 0.4555 & 1.4952 & 1.0864 & 0.7896 & 80.42\% \\
Epoch 1 (Best Model)    & 0.3666 & 0.5997 & 0.6753 & 0.4749 & 79.08\% \\
Epoch 2                 & 0.3053 & 0.5159 & 0.6419 & 0.4148 & 78.63\% \\
Epoch 3 (Convergence)   & 0.2598 & 0.5183 & 0.6412 & 0.3877 & 78.72\% \\
\midrule
\textbf{Net Relative Change} & \textbf{-42.9\%} & \textbf{-65.3\%} & \textbf{-41.0\%} & \textbf{-50.9\%} & \textbf{Preserved} \\
\bottomrule
\end{tabular}%
}
\end{table}

\section{Confidence Calibration \& Safety Gating}
\subsection{Temperature Scaling}
Deep classifiers are frequently uncalibrated, producing overconfident probabilities on clinical errors~\cite{guo2017calibration}. We perform post-hoc temperature scaling on a held-out calibration split comprising 401 patient records (2,406 item-question evaluations). The calibrated confidence distribution is defined as:
\begin{equation}
\hat{P}(Y = c \mid \mathbf{x}) = \frac{\exp(z_c / T)}{\sum_{j=1}^3 \exp(z_j / T)}
\end{equation}
Optimizing $T$ via negative log-likelihood minimization yields an optimal parameter of \textbf{$T = 1.1911$}. Temperature scaling reduced calibration NLL from 0.4791 to 0.4696 while maintaining a well-calibrated mean confidence of 83.66\%.

\subsection{Clinician Review Safety Trigger}
To prevent automated diagnostic failures, CEIS implements a deterministic safety gating rule flagging cases for manual clinician audit:
\begin{equation}
\begin{aligned}
\text{Flag}_{\text{review}}(\mathbf{x}) = \;&\left(\max_{c} \hat{P}_c < 0.70\right) \;\lor\; (\hat{y} = \text{contradicted}) \\
&\lor\; (\mathcal{S}_{\text{sev}} > 0.60)
\end{aligned}
\end{equation}
When triggered, the system surfaces the structured reasoning trail, calibrated confidence, and exact text spans to attending medical staff.

\section{Experimental Evaluation \& Head-to-Head Comparison}
To evaluate whether fine-tuning improved multi-question capability, we benchmarked the baseline model (\texttt{pavanmantha/mednli\_laya}) against the fine-tuned CEIS model (\texttt{reddyshailesh1/ceis}) across representative clinical challenge scenarios.

\begin{table*}[t]
\centering
\caption{Qualitative Head-to-Head Comparison: Baseline Semantic Failure vs. CEIS Multi-Head Alignment}
\label{tab:head_to_head}
\begin{tabular}{@{}p{3.8cm}p{3.2cm}p{3.2cm}p{6.4cm}@{}}
\toprule
\textbf{Clinical Scenario \& Input Pair} & \textbf{Baseline Checkpoint} & \textbf{CEIS (Fine-Tuned)} & \textbf{Clinical Impact \& Analysis} \\
\midrule
\textbf{Case 1: Corroborated Evidence}\newline
\textit{Claim:} Patient was diagnosed with pulmonary embolism.\newline
\textit{Evidence:} CT pulmonary angiogram demonstrated acute filling defect in the right main pulmonary artery consistent with pulmonary embolism. &
Verdict: \textbf{SUPPORTED} (0.958)\newline
Mentioned: 0.994\newline
Needs More: \textbf{0.989 (FAIL)} &
Verdict: \textbf{SUPPORTED} (0.956)\newline
Mentioned: 0.803\newline
Needs More: \textbf{0.346 (PASS)} &
\textbf{Severe Semantic Confusion in Baseline:} Baseline correctly detects entailment, but simultaneously outputs 98.9\% probability that more evidence is needed. CEIS eliminates this contradiction, recognizing CT angiography as definitive diagnostic proof. \\
\midrule
\textbf{Case 2: Explicit Clinical Rule-Out}\newline
\textit{Claim:} Patient has acute kidney injury.\newline
\textit{Evidence:} Serum creatinine was stable at 0.7\,mg/dL with robust urine output of 150\,mL/hr, ruling out acute kidney injury. &
Verdict: \textbf{CONTRADICTED} (0.964)\newline
Mentioned: \textbf{0.014 (FAIL)}\newline
Needs More: 0.057 &
Verdict: \textbf{CONTRADICTED} (0.909)\newline
Mentioned: \textbf{0.800 (PASS)}\newline
Needs More: 0.332 &
\textbf{Topic Recognition Failure in Baseline:} Baseline outputs 1.4\% for mentioned, falsely asserting that the note never discussed AKI despite explicit textual rule-out. CEIS accurately identifies explicit topic discussion (80.0\%). \\
\bottomrule
\end{tabular}
\end{table*}

\subsection{Head-to-Head Analysis}
As detailed in Table~\ref{tab:head_to_head}, the baseline checkpoint exhibits catastrophic semantic head incoherence when queried on non-verdict primitives:
\begin{itemize}
    \item \textbf{Case 1 (Definitive Imaging Proof):} The baseline predicts \texttt{needs\_more\_evidence = 0.989}. In a clinical workflow, this would falsely alert clinicians that diagnostic proof is missing when a definitive pulmonary embolism is established by CT angiography. CEIS resolves this, reducing the probability to 0.346.
    \item \textbf{Case 2 (Explicit Diagnostic Rule-Out):} The baseline predicts \texttt{explicitly\_mentioned = 0.014}. The baseline falsely concludes that acute kidney injury was never discussed. CEIS correctly predicts 0.800, successfully distinguishing between unmentioned claims and explicit rule-outs.
\end{itemize}

\subsection{Out-of-Domain Generalization on PubMedQA}
We evaluated zero-shot transfer on PubMedQA~\cite{jin2019pubmedqa} (1,000 biomedical abstract QA pairs mapped to supported/unsupported/contradicted). CEIS achieves zero-shot accuracy of 61.4\% and Macro-F1 of 0.582, confirming that non-autoregressive decision representations generalize across biomedical literature without catastrophic forgetting.

\section{Software Architecture \& Deployment}
The complete system is organized into modular tiers (Fig.~\ref{fig:arch}):
\begin{itemize}
    \item \textbf{Chunking \& Span Extraction (\texttt{inference/}):} For documents exceeding 512 tokens, a sliding-window chunker (400-token window, 50-token overlap) aggregates token representations. An evidence highlighter computes token-level salience, extracting supporting sentences in green and contradicting sentences in red.
    \item \textbf{Contradiction Matrix Scanner:} Evaluates all $N \times M$ statement pairs between disparate notes (e.g., Admission vs. Discharge), generating a visual contradiction heatmap.
    \item \textbf{FastAPI REST Service (\texttt{app/api.py}):} Delivers asynchronous sub-50\,ms endpoints (\texttt{/verify}, \texttt{/scan\_contradictions}, \texttt{/sufficiency}).
    \item \textbf{Streamlit Clinical Interface (\texttt{app/frontend/}):} Interactive clinical dashboard supporting multi-document verification, risk visualization, and one-click JSON/CSV audit logging.
\end{itemize}

\begin{figure}[t]
\centering
\fbox{
\begin{minipage}{0.92\columnwidth}
\small
\textbf{System Pipeline Flow:}\\[2pt]
1. Clinical Input $(C, E) \rightarrow$ Sliding-Window Chunker\\[2pt]
2. Joint Tokenization $\rightarrow$ Shared Non-Autoregressive Backbone\\[2pt]
3. Parallel Evaluation Heads:\\[1pt]
\hspace*{8pt}$\bullet$ Choice Head: Verdict $(P = 0.956)$\\[1pt]
\hspace*{8pt}$\bullet$ Noul Head: Explicitly Mentioned $(P = 0.803)$\\[1pt]
\hspace*{8pt}$\bullet$ Noul Head: Needs More Evidence $(P = 0.346)$\\[1pt]
\hspace*{8pt}$\bullet$ Score Head: Sufficiency Index $(S = 0.90)$\\[2pt]
4. Temperature Calibration ($T = 1.1911$) \& Safety Review Gating\\[2pt]
5. HTML Evidence Span Highlighting \& REST / UI Dispatch
\end{minipage}
}
\caption{Inference flow of the Clinical Evidence Intelligence System.}
\label{fig:arch}
\end{figure}

\section{Limitations \& Ethical Considerations}
Several limitations should be noted. First, the model is trained primarily on adult ICU documentation from MIMIC-III; performance on pediatric, outpatient, or non-English documentation remains to be evaluated. Second, silver boolean supervision was derived programmatically; prospective multi-annotator clinician labeling will further refine sufficiency scoring. Finally, CEIS is an assistive evidence verification tool designed to augment clinician efficiency; it does not replace physician clinical judgment.

\section{Conclusion \& Future Work}
We presented the Advanced Clinical Evidence Intelligence System (CEIS), extending LAYA's non-autoregressive architecture into a multi-task clinical decision framework. By evaluating entailment, explicit mention, sufficiency, and severity simultaneously, CEIS eliminates the hallucination risks of generative LLMs while resolving the semantic limitations of standard single-label classifiers. Fine-tuning on an NVIDIA RTX 6000 Ada (48\,GB) resolved acute semantic head failures, yielding calibrated, trustworthy predictions. Future work will investigate cross-lingual clinical transfer and real-time electronic health record integration.

\section*{Acknowledgment}
The authors acknowledge the open-source contributors of LAYA and the PhysioNet research community for providing access to the MIMIC-III and MedNLI datasets.

\bibliographystyle{IEEEtran}
\begin{thebibliography}{1}

\bibitem{romanov2018lessons}
A.~Romanov and C.~Shivade, ``Lessons from natural language inference in the clinical domain,'' in \emph{Proc. Conf. Empir. Methods Nat. Lang. Process. (EMNLP)}, Brussels, Belgium, 2018, pp. 1586--1596.

\bibitem{jin2019pubmedqa}
Q.~Jin, B.~Dhingra, Z.~Liu, W.~Cohen, and X.~Lu, ``{PubMedQA}: A dataset for biomedical research question answering,'' in \emph{Proc. Conf. Empir. Methods Nat. Lang. Process. (EMNLP)}, Hong Kong, 2019, pp. 2567--2577.

\bibitem{alsentzer2019publicly}
E.~Alsentzer, J.~Murphy, W.~Boag, W.-H. Weng, D.~Jeng, T.~Naumann, and M.~McDermott, ``Publicly available clinical {BERT} embeddings,'' in \emph{Proc. 2nd Clin. Nat. Lang. Process. Workshop}, Minneapolis, MN, 2019, pp. 72--78.

\bibitem{gu2021domain}
Y.~Gu, R.~Tinn, H.~Cheng, M.~Lucas, N.~Usuyama, X.~Zheng, P.~Gao, J.~Han, and H.~Poon, ``Domain-specific language model pretraining for biomedical natural language processing,'' \emph{ACM Trans. Comput. Healthcare}, vol.~3, no.~1, pp. 1--23, 2021.

\bibitem{laya2026}
N.~K.~M., ``{LAYA}: Non-autoregressive decision model architecture,'' 2026. [Online]. Available: \url{https://github.com/NandhaKishorM/laya}

\bibitem{guo2017calibration}
C.~Guo, G.~Pleiss, Y.~Sun, and K.~Q. Weinberger, ``On calibration of modern neural networks,'' in \emph{Proc. 34th Int. Conf. Mach. Learn. (ICML)}, Sydney, Australia, 2017, pp. 1321--1330.

\bibitem{johnson2016mimic}
A.~E. Johnson, T.~J. Pollard, L.~Shen, L.-W.~H. Lehman, M.~Feng, M.~Ghassemi, B.~Moody, P.~Szolovits, L.~A. Celi, and R.~G. Mark, ``{MIMIC-III}, a freely accessible critical care database,'' \emph{Sci. Data}, vol.~3, no.~1, p. 160035, 2016.

\bibitem{reddy2026ceis}
S.~Reddy and Karthikeya, ``Clinical Evidence Intelligence System ({CEIS}) on {LAYA},'' Hugging Face, 2026. [Online]. Available: \url{https://huggingface.co/reddyshailesh1/ceis}

\end{thebibliography}

\end{document}
